\chapter{Automated Reasoning}
\label{ch:automated-reasoning}

Automated reasoning separates a calculus from a search procedure.  A calculus
determines which inferences are licensed.  A search procedure decides which
licensed inference to try next.  Soundness is a property of the inferences;
termination and efficiency are properties of the search procedure.  The
handbook source for the calculi and procedures of this chapter is
\cite{RobinsonVoronkov2001}.

\section{Propositional satisfiability}
\label{sec:automated-reasoning-01}

The terms \emph{propositional literal}, \emph{clause}, and
\emph{conjunctive normal form} retain their definitions from the classical
propositional chapter.  In particular, a clause \(C\) is a finite set of
propositional literals, not a formula, and
\(\operatorname{form}_{\CPL}(C)\) is its exact formula interpretation.  Let
\(S\) be a finite set of clauses.  A valuation satisfies \(S\) iff it
satisfies \(\operatorname{form}_{\CPL}(C)\) for every \(C\in S\).  The empty
set of clauses is satisfied by every valuation, whereas a set containing the
empty clause is satisfied by none.

\begin{definition}[Unit propagation]
Let \(S\) be a finite set of clauses.  If \(S\) contains a unit clause
\(\set L\), \emph{unit propagation} assigns \(L\) true, deletes every clause
containing \(L\), and deletes the complementary literal from every remaining
clause.
\end{definition}

\begin{definition}[Pure-literal elimination]
If a propositional variable occurs in a finite clause set with only one
polarity, \emph{pure-literal elimination} assigns the truth value that
satisfies all of those occurrences and deletes the clauses containing them.
\end{definition}

The Davis--Putnam--Logemann--Loveland (DPLL) procedure recursively applies
unit
propagation, optionally applies pure-literal elimination, and otherwise branches
on an unassigned variable \cite{DavisPutnam1960,DavisLogemannLoveland1962}.
The following pseudocode specifies the logical core; it does not specify a
branching heuristic.  The set \(V\) of variables occurring in the original
input is retained unchanged by every recursive call.

\begin{lstlisting}[caption={Logical core of a DPLL satisfiability procedure}]
DPLL(S, valuation, V):
  repeatedly unit-propagate S, extending valuation by each forced value
  if S has no clauses:
    extend valuation to every unassigned variable in V with value false
    return SAT with the resulting total valuation
  if S contains the empty clause: return UNSAT
  choose an unassigned variable p
  r := DPLL(S with p=true, valuation extended by p=true, V)
  if r is SAT: return r
  return DPLL(S with p=false, valuation extended by p=false, V)
\end{lstlisting}

\begin{theorem}[Correctness of the displayed DPLL procedure]
For every finite set \(S\) of clauses, the displayed
procedure terminates.  It returns \(\SAT\) with a satisfying valuation iff
\(S\) is satisfiable, and it returns \(\UNSAT\) iff \(S\) is
unsatisfiable.
\end{theorem}
\begin{proof}
Maintain the invariant that every deleted original clause already contains a
literal made true by the current partial valuation, and every undeleted clause
is its residual after removing literals made false.  Unit propagation preserves
this invariant: every satisfying extension of a unit clause makes its literal
true.  If no residual clauses remain, every original clause is already true;
assigning arbitrary values to the remaining input variables cannot falsify a
clause that already contains a true literal.  Thus the displayed completion by
\(\False\) returns a total satisfying valuation on \(V\).

At a branch on \(p\), every total valuation makes either \(p\) true or false,
so the current residual set is satisfiable iff at least one branch is
satisfiable.  Each branch assigns a previously unassigned variable; there are
finitely many variables, so every path terminates.  Induction on the number of
unassigned variables now establishes both return specifications.
\end{proof}

